\subsection{\texorpdfstring{埃尔米特型与 \(\varepsilon\)-埃尔米特型。}{埃尔米特型与 epsilon-埃尔米特型。}}

定义 1. --- 设 \(\varepsilon\) 是 A 的中心的一个元素。称半双线性型 \(\Phi\)（E 上的）满足 \(\Phi(x, y) = \varepsilon \overline{\Phi(y, x)}\)（对 E 中一切 \(x\) 与 \(y\)）者为 \(\varepsilon\)-埃尔米特型。1-埃尔米特型（相应地 \((-1)\)-埃尔米特型）称为埃尔米特型（相应地反埃尔米特型）。

当 J 为恒等（这蕴含 A 交换）时，（关于 J 的）埃尔米特型（相应地反埃尔米特型）正是对称（相应地反对称）双线性型（第 III 章，§ 5，第 1 小节，定义 2）。回忆，交错双线性型（第 III 章，§ 5，第 2 小节，定义 4）是反对称的；反之，若在 A 中关系 \(2a = 0\) 蕴含 \(a = 0\)，则逆命题成立。

关于 ε-埃尔米特型的正交关系（§ 1，第 3 小节）显然是对称的（参见习题 1）。

若 \(\alpha\) 是 A 的可逆元素，则映射 T: \(\lambda \to \alpha^{-1}\overline{\lambda}\alpha\) 是 A 的一个反自同构，且容易验证型 \(\Phi\alpha\) 关于 T 是半线性的。此外，若 \(\alpha = \overline{\alpha}\)，则 T 是对合的，且若 \(\Phi\) 是 \(\varepsilon\)-埃尔米特型，则 \(\Phi\alpha\) 也是；事实上有

\[
(\lambda^ {\mathrm{T}}) ^ {\mathrm{T}} = \alpha^ {- 1} (\overline {{\alpha^ {- 1} \lambda \alpha}}) \alpha = \alpha^ {- 1} \overline {{\alpha}} \lambda \overline {{\alpha}} ^ {- 1} \alpha = \lambda
\]

\[
\Phi (y, x) \alpha = \varepsilon \overline {{\Phi (x , y)}} \alpha = \varepsilon (\Phi (x, y) \alpha) ^ {\mathrm{T}}.
\]

特别地，当 A 是一个体时，A 的中心的元素 \(\alpha\) 满足 \(\bar{\alpha} = \alpha\) 者组成 A 的一个子域 K，而 E 上（关于 J 的）\(\varepsilon\)-埃尔米特型组成 K 上的一个向量空间。

注. --- 1) 若 \(\Phi\) 是 E 上的 \(\varepsilon\)-埃尔米特型，则 \(\Phi(x, y) = \varepsilon \Phi(x, y) \bar{\varepsilon}\) 对 E 中一切 \(x, y\) 成立。因此，若 \(\Phi\) 取可逆值，则 \(\varepsilon \bar{\varepsilon} = 1\)。

\begin{enumerate}
\def\labelenumi{\arabic{enumi})}
\setcounter{enumi}{1}
\tightlist
\item
  若 A 的中心存在可逆元素 i 使得 \(\bar{i}=i\varepsilon\)，则要使 \(\Phi\) 是 \(\varepsilon\)-埃尔米特型，必须且只需 \(i\Phi\) 是埃尔米特型。
\end{enumerate}

由于映射 \((y, x) \to \overline{\Phi(x, y)}\) 关于 J 是半双线性的，要使 \(\Phi\) 是 \(\varepsilon\)-埃尔米特型，必须且只需有 \(\Phi(y, x) = \varepsilon \overline{\Phi(x, y)}\)（当 \(x\) 与 \(y\) 取遍 E 的一个生成组时）。特别地，若 E 具有有限基 \((e_i)_{1 \leqslant i \leqslant n}\)，则要使 E 上的半双线性型 \(\Phi\) 是 \(\varepsilon\)-埃尔米特型，必须且只需它的矩阵 \(R = (\rho_{ij}) = (\Phi(e_i, e_j))\) 满足关系 \(\rho_{ji} = \varepsilon \overline{\rho_{ij}}\)（对一切 \(i, j\)），即 \(^t R = \varepsilon \overline{R}\)；具有这一性质的矩阵 \(R\) 称为 \(\varepsilon\)-埃尔米特矩阵。当 \(\varepsilon = 1\)（相应地 -1）时，称 \(R\) 关于反自同构 J 是埃尔米特（相应地反埃尔米特）矩阵。当 J 为恒等（从而 A 交换）时，埃尔米特（相应地反埃尔米特）矩阵 \(R\) 满足 \(^t R = R\)（相应地 \(^t R = -R\)）；这时称 \(R\) 为对称（相应地反对称）矩阵。要使 \(\Phi\) 是交错型，必须且只需它的矩阵是反对称的，且此外 R 的对角项全为零；具有这些性质的矩阵称为交错矩阵。

设 \(\Phi\) 是 E 上的一个半双线性型，\(s_{\Phi}\) 与 \(d_{\Phi}\) 是 \(\Phi\) 左相伴与右相伴的 E 到 E* 中的映射（§ 1，第 6 小节）。要使 \(\Phi\) 是 \(\varepsilon\)-埃尔米特型，必须且只需 \(\langle x, s_{\Phi}(y) \rangle = \bar{\varepsilon} \langle x, d_{\Phi}(y) \rangle\) 对 E 的一切元素 \(x, y\) 成立，即 \(s_{\Phi} = \bar{\varepsilon} d_{\Phi}\)，或者等价地，\(\langle x, d_{\Phi}(y) \rangle = \varepsilon \langle x, s_{\Phi}(y) \rangle\)，即 \(d_{\Phi} = \varepsilon s_{\Phi}\)。

设 \(\Phi\) 是一个 \(\varepsilon\)-埃尔米特型，使得映射 \(d_{\Phi}\) （E 到 E* 中的、右相伴于 \(\Phi\) 的）是双射。于是对 E 的每个自同态 \(u\) 有

\[
u ^ {* *} = \varepsilon \bar {\varepsilon} u.\tag{1}
\]

事实上，当 \(x\) 与 \(y\) 是 \(\mathbf{E}\) 中的任意元素时，有

\[
\begin{array}{r l} \Phi (x, u ^ {* *} (y)) & = \Phi (u ^ {*} (x), y) = \varepsilon \overline {{\Phi (y , u ^ {*} (x))}} = \varepsilon \overline {{\Phi (u (y) , x)}} \\ & = \varepsilon \Phi (x, u (y)) \bar {\varepsilon} = \Phi (x, \varepsilon \bar {\varepsilon} u (y)) \end{array}
\]

因此 \(u^{**}(x) = \varepsilon \bar{\varepsilon} u(x)\)，因为 \(\Phi\) 非退化。

若 \(\Phi\) 是 \(\varepsilon\)-埃尔米特型，使得映射 \(s_{\Phi}\) 与 \(d_{\Phi}\) 均为双射，则逆型 \(\hat{\Phi}\)（\(\Phi\) 的逆型）（\(\S 1\)，第 7 小节）是 \(\bar{\varepsilon}\)-埃尔米特型。事实上，为简便起见令 \(s = s_{\Phi}\)、\(d = d_{\Phi}\)，由 \(d = \varepsilon s\) 可得 \(s^{-1} = \bar{\varepsilon} d^{-1}\)（\(s\) 是半线性的）。于是，对 E 中任意的 \(u\)、\(\nu\)，有

\[
\hat {\Phi} (u, \nu) = \Phi (s ^ {- 1} (u), d ^ {- 1} (\nu)) = \bar {\varepsilon} \Phi (d ^ {- 1} (u), d ^ {- 1} (\nu)),
\]

从而

\[
\widehat {\Phi} (\nu , u) = \overline {{\varepsilon \varepsilon}} \overline {{\Phi (d ^ {- 1} (u) , d ^ {- 1} (\nu))}} = \overline {{\varepsilon}} \widehat {\Phi} (u, \nu),
\]

因为 ε 属于 A 的中心。

最后，当环 A 交换时，\(\varepsilon\)-埃尔米特型 \(\Phi\) 到 E 的张量幂与外幂 \(\bigotimes^{p}\mathbf{E}\) 和 \(\bigwedge^{p}\mathbf{E}\) 的典范扩张是 \(\varepsilon^{p}\)-埃尔米特型，这由 § 1 第 9 小节公式 (35) 与 (37) 立即得出。

